% ===================================================================
%  TABLEAU N° 3 — Enfance et Jeunesse.
%  Feuillets 19 a 24 du fac-simile = folios 17 a 22.
%  Correspondance : folio imprime = numero de feuillet - 2.
%
%  Ca tableau (kom lo montras texto/io/12-tabelo-03.tex) kombinas du
%  ceni familiala : « Un bapteme au village » e « Une fete publique ».
%  La numerado d'alineas ripetas de 1 a singla ceno (voyez CONSIGNE-
%  RELEVE.md § 6) : la cles %%K restas do t03-01-1 ... t03-24-1 por la
%  unesma ceno, pose denove t03-01-1 ... t03-16-1 por la duesma — la
%  distinction c1/c2 sera enkondukita aparte da outils/ceni.py.
%
%  Le francais ne conserve que DEUX rangs de titre (« Premiere scene. »
%  et « Un bapteme au village. », puis « Deuxieme scene. » et « Une
%  fete publique. ») la ou l'ido en pose quatre par scene (« La Kirko.»,
%  « La frukto-gardeno. », « La baptal defilo. », « La spektanti. »,
%  « Navigal ludi e regati. », « La kuro. », « La ferio. ») : ces
%  cles-la restent donc sans vis-a-vis, comme dans texto/fr/10-tableau-
%  01.tex. C'est un ecart REEL entre les deux editions.
%
%  L'edition francaise s'acheve UNE PAGE PLUS TOT que l'ido : le
%  tableau 3 tient ici en 6 feuillets (19-24, folios 17-22) contre 7
%  chez Rochelle (19-25, folios 17-23) ; le dernier alinea francais (16.
%  Le dompteur...) correspond a l'alinea 16 ido, mais l'ido continue
%  encore avec un alinea 15/16 different — verifier au controle final.
%
%  ATTENTION — LE FAC-SIMILE FRANCAIS EST UNE PRISE DE VUE, contraste
%  inegal d'une page a l'autre : plusieurs renvois ont du etre regeneres
%  a 2600 px pour etre lus avec certitude (ex. « rûche » p.18, coquille
%  imprimee pour « ruche », conservee telle quelle).
% ===================================================================

% --- Feuillet 19 (folio 17 : apertajo di la tabelo) -----------------
\begin{VUpage}[19]{}
%%K t03-apar-1 apar
\VUsaut{3.0mm}
\VUtitre{15.0pt}{60}{TABLEAU N\textsuperscript{o} 3}
\VUsaut{1.6mm}
\VUtitre{11.4pt}{0}{Enfance et Jeunesse.}
\VUsaut{3.4mm}

%%K t03-apar-2 apar
\VUblancAlinea
(Les tableaux n\textsuperscript{os} 3 et 4 sont combinés de façon à\nl
présenter en quatre \VUgras{scènes familiales} les notions\nl
essentielles relatives au \VUgras{temps}, à l'\VUgras{âge}, à la\nl
\VUgras{famille}, aux \VUgras{vêtements} et à la \VUgras{santé.})

\VUsaut{4.0mm}
%%K t03-tit-1 sub
\VUcentre{12.0pt}{0}{Première scène.}
\VUsaut{3.0mm}

%%K t03-tit-2 sub
\VUpk{10.2pt}{40}{Un baptême au village.}
\VUsaut{3.0mm}

%%K t03-c1-01-1 p
\VUblancAlinea
1. --- Nous sommes au \VUgras{printemps}, au mois de\nl
\VUgras{mars}, \VUgras{avril} ou \VUgras{mai}, un jour de la \VUgras{semaine},\nl
\VUgras{lundi}, \VUgras{mardi} ou \VUgras{mercredi}. Il est \VUgras{dix heures} du\nl
\VUgras{matin}.

%%K t03-c1-02-1 p
\VUblancAlinea
2. --- Le \VUgras{temps} est beau, l'\VUgras{air} pur. Le soleil\nl
brille. Quelques \VUgras{nuages} \textsuperscript{(2)} légers montent dans le\nl
\VUgras{ciel} \textsuperscript{(1)} bleu.

%%K t03-c1-03-1 p
\VUblancAlinea
3. --- Les \VUgras{arbres} ont à peine des \VUgras{feuilles}. Les\nl
minces \VUgras{peupliers} \textsuperscript{(3)} et le grand \VUgras{platane} \textsuperscript{(17)} ne sont\nl
encore couverts que de bourgeons.

%%K t03-c1-04-1 p
\VUblancAlinea
4. --- Les \VUgras{hirondelles} \textsuperscript{(4)} reviennent et voltigent\nl
en poussant de joyeux cris autour de la \VUgras{flèche} \textsuperscript{(7)} du\nl
\VUgras{clocher} \textsuperscript{(5)} qui surmonte l'\VUgras{église} \textsuperscript{(6)} du village.

%%K t03-c1-05-1 p
\VUblancAlinea
5. --- Elle est bien simple cette église avec son\nl
grand \VUgras{porche} \textsuperscript{(13)} et sa grosse \VUgras{horloge} \textsuperscript{(8)}. La\nl
\VUgras{petite aiguille} \textsuperscript{(9)} est sur le chiffre X; elle marque\nl
les \VUgras{heures}. La \VUgras{grande} \textsuperscript{(10)} est sur midi; elle mar\cc
que les \VUgras{minutes}.

%%K t03-c1-06-1 p
\VUblancAlinea
6. --- Dans le clocher, les \VUgras{cloches} \textsuperscript{(11)} sonnent\nl
joyeusement. A l'extrémité du toit se dresse une\nl
petite \VUgras{croix} \textsuperscript{(12)} de pierre.

%%K t03-c1-07-1 p
\VUblancAlinea
7. --- L'église n'a pas seulement un grand \VUgras{por}\cc
\VUgras{tail} \textsuperscript{(13)}, elle a aussi une petite porte latérale. Par
\parplein
\end{VUpage}

% --- Feuillet 20 (folio 18) ------------------------------------------
\begin{VUpage}[20]{18}
%%K t03-c1-07-1 p suite
\VUcontinue
cette porte M. le \VUgras{Curé} \textsuperscript{(15)}, vêtu de sa \VUgras{soutane}, sort\nl
de la \VUgras{sacristie} \textsuperscript{(14)} et se rend au \VUgras{presbytère} \textsuperscript{(19)}.

%%K t03-c1-08-1 p
\VUblancAlinea
8. --- Près de lui, voici la \VUgras{laitière} \textsuperscript{(24)} avec son\nl
\VUgras{âne} \textsuperscript{(23)}. Elle revient de porter à la ville son bon lait\nl
pour les petits enfants.

%%K t03-c1-09-1 p
\VUblancAlinea
9. --- Maître Aliboron est tout content de cette\nl
course matinale. Il aspire avec délices l'air embaumé\nl
par le parfum des \VUgras{fleurs} \textsuperscript{(20)} qui couvrent les \VUgras{arbres}\nl
\VUgras{fruitiers} \textsuperscript{(18)} du \VUgras{verger} voisin. Ils sont tout roses\nl
et tout blancs ces \VUgras{pêchers} \textsuperscript{(18)} et ces \VUgras{abricotiers} \textsuperscript{(19)},\nl
et promettent une belle récolte.

%%K t03-c1-10-1 p
\VUblancAlinea
10. --- En attendant, les petites \VUgras{abeilles} \textsuperscript{(22)} de la\nl
\VUgras{rûche} \textsuperscript{(24)} vont y butiner leur doux \VUgras{miel} en bour\cc
donnant.

%%K t03-c1-11-1 p
\VUblancAlinea
11. --- Mais que se passe-t-il ? Voilà les enfants du\nl
village qui accourent et font fuir les \VUgras{oies} \textsuperscript{(25)}\nl
effrayées. C'est le \VUgras{baptême} du fils du notaire qui\nl
sort de l'église.

%%K t03-c1-12-1 p
\VUblancAlinea
12. --- Le petit Jacques, dans son \VUgras{berceau} \textsuperscript{(26)}, se\nl
réveille au joyeux son des cloches, et sa sœur Made\cc
leine, qui veut voir, elle aussi, le cortège, demande à\nl
sa mère de venir la peigner et l'habiller sur le seuil\nl
de la maison.

%%K t03-c1-13-1 p
\VUblancAlinea
13. --- La mère y consent. Elle met son \VUgras{foulard} \textsuperscript{(28)},\nl
son \VUgras{corsage} \textsuperscript{(29)} et son \VUgras{tablier} \textsuperscript{(30)}, et elle vient\nl
s'asseoir sur une petite chaise devant la porte. Made\cc
leine n'a que son \VUgras{jupon} \textsuperscript{(31)} et son \VUgras{corset} \textsuperscript{(32)}.

%%K t03-c1-14-1 p
\VUblancAlinea
14. --- Comme elle est heureuse ! Elle oublie ses\nl
fleurs favorites, ses \VUgras{œillets} \textsuperscript{(34)}, où se posent les\nl
\VUgras{papillons} \textsuperscript{(35)}, ses \VUgras{tulipes} \textsuperscript{(36)}, ses \VUgras{muguets} \textsuperscript{(37)},\nl
dans des \VUgras{pots} \textsuperscript{(38)}, sur le \VUgras{mur} \textsuperscript{(33)}; et ses \VUgras{roses} \textsuperscript{(39)},\nl
qui forment une si belle haie. Elle contemple les\nl
belles toilettes des dames du cortège.

%%K t03-c1-15-1 p
\VUblancAlinea
15. --- Les garçons du village ne s'occupent guère\nl
des toilettes. Ils se précipitent et se bousculent pour
\parplein
\end{VUpage}

% --- Feuillet 21 (folio 19) ------------------------------------------
\begin{VUpage}[21]{19}
%%K t03-c1-15-1 p suite
\VUcontinue
avoir des \VUgras{bonbons} \textsuperscript{(55)} ou des \VUgras{sous} \textsuperscript{(52)}. L'un d'eux\nl
pleure \textsuperscript{(40)}.

%%K t03-c1-16-1 p
\VUblancAlinea
16. --- L'autre a marché sur la patte du \VUgras{petit}\nl
\VUgras{chien} \textsuperscript{(42)}, qui se sauve en criant et met en fuite la\nl
\VUgras{poule} \textsuperscript{(43)} et ses \VUgras{poussins} \textsuperscript{(44)}.

%%K t03-c1-17-1 p
\VUblancAlinea
17. --- En tête du cortège marche la \VUgras{nourrice} \textsuperscript{(45)}.\nl
Elle a un beau \VUgras{bonnet} \textsuperscript{(16)} avec de longs rubans.\nl
Elle porte le \VUgras{nouveau-né} \textsuperscript{(47)}, tout habillé de \VUgras{den}\cc
\VUgras{telles} \textsuperscript{(48)}.

%%K t03-c1-18-1 p
\VUblancAlinea
18. --- Derrière la nourrice viennent le \VUgras{parrain} \textsuperscript{(49)}\nl
et la \VUgras{marraine} \textsuperscript{(53)}. Ils distribuent les bonbons et\nl
les sous. Le parrain a des \VUgras{gants} \textsuperscript{(51)} et un \VUgras{chapeau}\nl
\VUgras{haut de forme} \textsuperscript{(50)}. La marraine tient à la main un\nl
joli \VUgras{bouquet} \textsuperscript{(54)}.

%%K t03-c1-19-1 p
\VUblancAlinea
19. --- Derrière eux nous apercevons l'\VUgras{oncle} \textsuperscript{(56)} et\nl
la \VUgras{tante} \textsuperscript{(58)} de l'enfant. La tante a un élégant \VUgras{cha}\cc
\VUgras{peau} \textsuperscript{(59)}, l'oncle une \VUgras{redingote} \textsuperscript{(57)} noire et un gilet\nl
blanc. Voici ensuite le \VUgras{cousin} \textsuperscript{(60)} et la \VUgras{cousine} \textsuperscript{(61)}.\nl
Celle-ci tient par la main la \VUgras{sœur} \textsuperscript{(62)} du nouveau-né,\nl
la petite Berthe.

%%K t03-c1-20-1 p
\VUblancAlinea
20. --- L'autre \VUgras{frère}\textsuperscript{(63)} de Berthe marche derrière.\nl
Il donne la main à une \VUgras{parente} \textsuperscript{(64)}. Au fond vien\cc
nent les \VUgras{amis} \textsuperscript{(65)} de la famille.

%%K t03-c1-21-1 p
\VUblancAlinea
21. --- Près du cortège, voici un \VUgras{mendiant} \textsuperscript{(66)}\nl
appuyé sur sa \VUgras{béquille} \textsuperscript{(67)}. Il demande l'aumône.

%%K t03-c1-22-1 p
\VUblancAlinea
22. --- De l'autre côté est un petit garçon très\nl
\VUgras{poli} \textsuperscript{(68)}. Il salue et dit « \VUgras{bonjour} » aux personnes\nl
qu'il connaît.

%%K t03-c1-23-1 p
\VUblancAlinea
23. --- Les gens du village sont sortis de leurs\nl
maisons pour voir passer le baptême. Nous voyons\nl
un \VUgras{paysan} \textsuperscript{(69)} avec son \VUgras{bonnet de coton} et près\nl
de lui une \VUgras{paysanne} \textsuperscript{(70)}. Puis une \VUgras{vieille femme} \textsuperscript{(71)},\nl
appuyée sur son \VUgras{bâton} \textsuperscript{(72)}. Enfin le \VUgras{meunier} \textsuperscript{(73)}\nl
avec son grand \VUgras{chapeau de feutre} \textsuperscript{(74)} et une jeune
\parplein
\end{VUpage}

% --- Feuillet 22 (folio 20) ------------------------------------------
\begin{VUpage}[22]{20}
%%K t03-c1-23-1 p suite
\VUcontinue
\VUgras{paysanne} chaussée de \VUgras{sabots} \textsuperscript{(75)}, qui apprend à mar\cc
cher à son petit enfant.

%%K t03-c1-24-1 p
\VUblancAlinea
24. --- Cette scène est très amusante.

\VUsaut{4.0mm}
%%K t03-tit-3 sub
\VUcentre{12.0pt}{0}{Deuxième scène.}
\VUsaut{3.0mm}

%%K t03-tit-4 sub
\VUpk{10.2pt}{40}{Une fête publique.}
\VUsaut{3.0mm}

%%K t03-c2-01-1 p
\VUblancAlinea
1. --- L'\VUgras{été} est venu. Le brûlant soleil du mois de\nl
\VUgras{juin} (\VUgras{juillet} ou \VUgras{août}) invite les jeunes gens à se\nl
baigner et à faire du canotage.

%%K t03-c2-02-1 p
\VUblancAlinea
2. --- Sur la rive droite du fleuve, les canotiers se\nl
déshabillent pour prendre part aux jeux nautiques et\nl
aux régates.

%%K t03-c2-03-1 p
\VUblancAlinea
3. --- L'un d'eux est en train de quitter ses \VUgras{sou}\cc
\VUgras{liers} \textsuperscript{(1)} ; il les délace (déboutonne). Ses deux cama\cc
rades sont déjà prêts. Ils n'ont plus que leur \VUgras{tricot} \textsuperscript{(2)}\nl
et leur \VUgras{caleçon} \textsuperscript{(3)}. Ils causent en attendant leurs\nl
amis. Ils parlent de la foire et de la fête qui a lieu\nl
sur l'autre rive.

%%K t03-c2-04-1 p
\VUblancAlinea
4. --- Par terre, auprès d'eux, sont étendus les\nl
\VUgras{vêtements} de leurs camarades : une \VUgras{paire} de \VUgras{bot}\cc
\VUgras{tines} \textsuperscript{(4)}, des \VUgras{souliers} \textsuperscript{(5)}, des \VUgras{chaussettes} \textsuperscript{(6)}, des\nl
\VUgras{chapeaux} de \VUgras{feutre} \textsuperscript{(7)} et de \VUgras{paille} \textsuperscript{(8)} et une \VUgras{cra}\cc
\VUgras{vate} \textsuperscript{(9)}.

%%K t03-c2-05-1 p
\VUblancAlinea
5. --- L'un des jeunes gens a soigneusement plié\nl
sa \VUgras{veste} \textsuperscript{(10)}. Il la pose sur une serviette, dans\nl
laquelle son voisin va serrer leurs deux \VUgras{costumes}.

%%K t03-c2-06-1 p
\VUblancAlinea
6. --- Un sixième garçon est en retard. Il a encore\nl
sa \VUgras{casquette} \textsuperscript{(11)} sur la tête. Il n'a quitté ni sa \VUgras{che}\cc
\VUgras{mise} \textsuperscript{(12)}, ni son \VUgras{col} \textsuperscript{(13)}, ni ses \VUgras{manchettes} \textsuperscript{(14)}. Il\nl
tient son \VUgras{gilet} \textsuperscript{(17)} à la main et il a encore son \VUgras{pan}\cc
\VUgras{talon} \textsuperscript{(16)} et ses \VUgras{bretelles} \textsuperscript{(15)}. Jamais il ne sera prêt\nl
pour la prochaine course.

%%K t03-c2-07-1 p
\VUblancAlinea
7. --- Près des jeunes gens, un sentier bordé de\nl
\VUgras{chardons} \textsuperscript{(18)} conduit au bord de la rivière, où le père
\parplein
\end{VUpage}

% --- Feuillet 23 (folio 21) ------------------------------------------
\begin{VUpage}[23]{21}
%%K t03-c2-07-1 p suite
\VUcontinue
Jacques prépare, dans les \VUgras{roseaux} \textsuperscript{(19)}, le \VUgras{feu}\nl
\VUgras{d'artifice} \textsuperscript{(20)} qui sera tiré le soir.

%%K t03-c2-08-1 p
\VUblancAlinea
8. Sur la rivière, plusieurs dames, s'abritant sous\nl
leurs ombrelles, se promènent en \VUgras{barque} \textsuperscript{(21)}. Elles\nl
suivent la course, qui est très intéressante. Dans chaque\nl
\VUgras{yole} \textsuperscript{(22)}, quatre \VUgras{rameurs} \textsuperscript{(24)} rament de toutes leurs\nl
forces, tandis que le \VUgras{barreur} \textsuperscript{(26)} tient le \VUgras{gouver}\cc
\VUgras{nail} \textsuperscript{(25)} et dirige la frêle embarcation. Les \VUgras{rames} \textsuperscript{(23)}\nl
frappent l'eau en cadence et les légers bateaux filent\nl
à une allure vertigineuse.

%%K t03-c2-09-1 p
\VUblancAlinea
9. --- Quelques jeunes garçons se disputent, près de\nl
la pile du pont, le prix du \VUgras{mât de beaupré}\textsuperscript{(28)}. L'un\nl
d'eux s'avance avec précaution sur le mât élastique\nl
et glissant pour attraper le petit \VUgras{drapeau} \textsuperscript{(29)} fixé à\nl
l'extrémité. Son \VUgras{concurrent} \textsuperscript{(30)} malheureux est tombé\nl
à l'eau. Il nage pour regagner la rive.

%%K t03-c2-10-1 p
\VUblancAlinea
10. --- Des jeunes gens plus âgés \VUgras{joutent} \textsuperscript{(32)} près\nl
du \VUgras{quai} \textsuperscript{(33)}. Un \VUgras{public} \textsuperscript{(31)} nombreux assiste à tous\nl
ces jeux, appuyé sur le \VUgras{parapet} du \VUgras{pont}\textsuperscript{(27)} et massé\nl
sur la \VUgras{rive}. Quelques spectateurs, cependant, se\nl
retournent pour suivre le jeu du \VUgras{mât de cocagne} \textsuperscript{(45)}\nl
ou pour admirer le \VUgras{cortège municipal} qui se rend\nl
à la distribution des prix.

%%K t03-c2-11-1 p
\VUblancAlinea
11. --- En tête, voici quelques \VUgras{gamins} \textsuperscript{(35)} qui pré\cc
\VUgras{cèdent} la \VUgras{musique} \textsuperscript{(36)} ; puis viennent le \VUgras{maire} \textsuperscript{(37)},\nl
les \VUgras{adjoints} et le \VUgras{conseil municipal} de la petite\nl
ville. Les \VUgras{pompiers} \textsuperscript{(38)} ferment la marche.

%%K t03-c2-12-1 p
\VUblancAlinea
12. --- Il y a beaucoup de monde sur le \VUgras{champ de}\nl
\VUgras{foire} \textsuperscript{(39)}. On vient voir les différentes \VUgras{baraques} :\nl
les \VUgras{chevaux de bois} \textsuperscript{(40)}, le \VUgras{cirque} \textsuperscript{(41)}, dont la\nl
représentation est déjà commencée ; le \VUgras{bal public} \textsuperscript{(42)},\nl
où jeunes gens et jeunes filles du pays se livrent au\nl
plaisir de la \VUgras{danse}.

%%K t03-c2-13-1 p
\VUblancAlinea
13. --- Dans le fond, nous apercevons les \VUgras{arènes} \textsuperscript{(44)},\nl
dans lesquelles le vaillant \VUgras{matador} espagnol combat\nl
le \VUgras{taureau} \textsuperscript{(45)} furieux ; puis les \VUgras{montagnes}
\parplein
\end{VUpage}

% --- Feuillet 24 (folio 22) ------------------------------------------
\begin{VUpage}[24]{22}
%%K t03-c2-13-1 p suite
\VUcontinue
\VUgras{russes} \textsuperscript{(49)}, le \VUgras{tir} \textsuperscript{(46)}, où l'on s'exerce à manier les\nl
\VUgras{armes à feu} et à tirer à la \VUgras{cible}.

%%K t03-c2-14-1 p
\VUblancAlinea
14. --- Près de là, voici le \VUgras{théâtre} \textsuperscript{(47)} \VUgras{forain},\nl
dans lequel on joue la \VUgras{comédie} et le \VUgras{drame}. Tout\nl
auprès, se trouve la baraque du \VUgras{géant} \textsuperscript{(48)} et du\nl
\VUgras{nain}. La \VUgras{taille} du premier est de 2\textsuperscript{m} 15 ; le second\nl
est à peine aussi grand qu'un enfant.

%%K t03-c2-15-1 p
\VUblancAlinea
15. --- Enfin, voici la grande \VUgras{ménagerie} \textsuperscript{(50)} avec\nl
ses \VUgras{animaux sauvages}, son \VUgras{tigre} \textsuperscript{(51)} aux \VUgras{griffes}\nl
puissantes, son \VUgras{lion} \textsuperscript{(52)} à la \VUgras{crinière} touffue, ses\nl
deux \VUgras{girafes} \textsuperscript{(53)} et son \VUgras{éléphant} \textsuperscript{(54)} qui se sert\nl
si adroitement de sa \VUgras{trompe}.

%%K t03-c2-16-1 p
\VUblancAlinea
16. --- Le \VUgras{dompteur} \textsuperscript{(55)} qui dresse ces animaux\nl
est sur la porte de la baraque.

\VUsaut{6.0mm}
\VUfilet{22mm}
\end{VUpage}
